\section*{第 \ref{ch:g7:stats} 章 --- 数据整理}

\begin{solution}{exo:g7:stats:1}
统计每个得分的频数：
\begin{center}
\renewcommand{\arraystretch}{1.2}
\begin{tabular}{l|ccccc|c}
得分 & $9$ & $11$ & $12$ & $14$ & $15$ & 合计 \\
\hline
频数 & $4$ & $3$ & $7$ & $3$ & $3$ & $20$ \\
\end{tabular}
\end{center}
\end{solution}

\begin{solution}{exo:g7:stats:2}
频率：$\frac{4}{20} = 20\,\%$，$\frac{3}{20} = 15\,\%$，
$\frac{7}{20} = 35\,\%$，$15\,\%$，$15\,\%$。和：
$20 + 15 + 35 + 15 + 15 = 100\,\%$。\checkmark
\end{solution}

\begin{solution}{exo:g7:stats:3}
第一次调查：$\frac{30}{50} = 0.6 = 60\,\%$。第二次：
$\frac{45}{90} = 0.5 = 50\,\%$。第一次调查的比例更高，
虽然其拥有一辆车的家庭绝对频数
较少。
\end{solution}

\begin{solution}{exo:g7:stats:4}
$\intco{5}{15}$：$8, 12, 9, 14, 11$ --- 频数为 $5$。
$\intco{15}{25}$：$23, 22, 18, 16$ --- 频数为 $4$。
$\intco{25}{35}$：$31, 27, 29, 25, 33$ --- 频数为 $5$。
$\intco{35}{45}$：$35$ --- 频数为 $1$。合计 $15$。\checkmark
\end{solution}

\begin{solution}{exo:g7:stats:5}
四个条形的高度分别为 $8$、$12$、$6$、$4$，
对应取值 $0$、$1$、$2$、$3$（刻度每隔 $2$ 较为合适）。
\end{solution}

\begin{solution}{exo:g7:stats:6}
共 $30$ 个家庭。频率：$\frac{8}{30}$，$\frac{12}{30}$，
$\frac{6}{30}$，$\frac{4}{30}$。角度（$\times 360^\circ$）：$96^\circ$，
$144^\circ$，$72^\circ$，$48^\circ$；和
为 $96 + 144 + 72 + 48 = 360^\circ$。\checkmark
\end{solution}

\begin{solution}{exo:g7:stats:7}
百分比：夏季 $\frac{162}{360} = 45\,\%$；春季
$\frac{90}{360} = 25\,\%$；秋季和冬季
各为 $\frac{54}{360} = 15\,\%$。在 $60$ 人中，选择夏季的有
$0.45 \times 60 = 27$ 人。
\end{solution}

\begin{solution}{exo:g7:stats:8}
A 班：$\frac{12}{30} = 40\,\%$。B 班：$\frac{14}{40} = 35\,\%$。
A 班步行者的比例更高，尽管
B 班步行的频数更多。
\end{solution}

\begin{solution}{exo:g7:stats:9}
纵轴从 $9\,700$ 开始时，三个条形的可见
高度分别为 $100$、$300$ 和 $400$ 个单位：最后一个看起来是
第一个的四倍，给人以销售额翻了两番的印象。若纵轴从 $0$ 开始，条形
几乎相等（$9\,800$ 到 $10\,100$ 约增长 $3\,\%$）：
真实的图景显示销售额几乎持平。
\end{solution}

\begin{solution}{exo:g7:stats:10}
频率之和为 $1$：$f = 1 - 0.35 - 0.4 = 0.25$。频数（总数
为 $80$）：$0.35 \times 80 = 28$；$0.4 \times 80 = 32$；
$0.25 \times 80 = 20$。检验：$28 + 32 + 20 = 80$。
\end{solution}

\begin{solution}{exo:g7:stats:11}
女生：$55\,\%$ 的 $400 = 220$；男生：$180$。在食堂就餐的女生：
$0.4 \times 220 = 88$；男生：$0.6 \times 180 = 108$。合计：
$88 + 108 = 196$ 名学生，即 $\frac{196}{400} = 49\,\%$ 的
学校 --- 介于 $40\,\%$ 与 $60\,\%$ 之间，更接近女生的比例，
因为女生人数更多。
\end{solution}

\begin{solution}{pb:g7:stats:1}
\textbf{1.} 频数：$120 > 90$，A 校获胜。频率：
$\frac{120}{400} = 0.30 = 30\,\%$，对比
$\frac{90}{250} = 0.36 = 36\,\%$：B 校获胜。奖项
应依据频率确定 --- B 校回收所用物品的占比更大，
A 校只是果汁喝得更多。

\textbf{2.} 教师：$\frac{18}{30} = 0.60 = 60\,\%$。学生：
$\frac{210}{600} = 0.35 = 35\,\%$。合计：
$\frac{18 + 210}{630} = \frac{228}{630} \approx 0.36 =
36\,\%$。学生人数是教师的二十倍，
因此合计频率几乎完全偏向学生一方 ---
人数多的一方起主导作用。

\textbf{3.} 频率之和必须为 $100\,\%$
（见\cref{def:g7:stats:frequency}），但
$45 + 30 + 20 + 10 = 105$：至少有一个扇形有误。

\textbf{4.} 频率为 $0.15$ 的频数：
$0.15 \times 40 = 6$。目前的频数为 $14 + 10 + 6 = 30$，因此
最后一项的频数为 $40 - 30 = 10$，其频率为
$\frac{10}{40} = 0.25$。（检验：$0.35 + 0.25 + 0.15 + 0.25 =
1$。）

\textbf{5.} 角度：$0.40 \times 360 = 144^\circ$，
$0.35 \times 360 = 126^\circ$，$0.25 \times 360 = 90^\circ$；
且 $144 + 126 + 90 = 360^\circ$，恰好构成一个圆。

\textbf{6.} $\frac{2.4}{10} = 0.24 = 24\,\%$ 的选票
已收回。

\textbf{7.} 电话簿和汽车登记册所列的多为富裕家庭，其投票
倾向与全国不同；$2.4$ 百万份答复是一次庞大的计数，
但抽取自错误的总体。与第 1 题相同：
重要的不是数了多少，而是数的是谁——
基于失真样本计算的频率，
描述的是样本，而不是全国。

\textbf{8.} 真实支持率：
$0.7 \times 1\,000 + 0.3 \times 9\,000 = 700 + 2\,700 = 3\,400$
$10\,000$ 人中支持者占 $34\,\%$。民意调查：
$0.7 \times 500 + 0.3 \times 500 = 350 + 150 = 500$ 名支持者，
即 $1\,000$ 人中占 $50\,\%$——高出十六个百分点，因为
富裕者占样本的一半，却只占全镇的十分之一。

\textbf{9.} 回复：$0.4 \times 200 = 80$ 人不满意，
$0.1 \times 800 = 80$ 人满意，因此共有 $160$ 份回复，其中
$80$ 人不满意。测得的不满意率为 $\frac{80}{160} = 50\,\%$。真实
不满意率为 $\frac{200}{1\,000} = 20\,\%$。没有人撒谎——
只是不满意者更愿意回复，于是调查更多地
听到了他们的声音。

\textbf{10.} 真实比例：富裕者占十分之一，故有 $100$
名富裕者与 $900$ 名普通受访者：
$0.7 \times 100 + 0.3 \times 900 = 70 + 270 = 340$ 人支持，
即 $1\,000$ 人中，修正后的调查预测为 $34\,\%$——与第 8 题的
真实值相同。五万份精心选取的访问，胜过两百万份
选择不当的访问。

\textbf{11.} 医院 A：轻症 $\frac{90}{100} = 90\,\%$；重症
$\frac{280}{400} = 70\,\%$；总体
为 $\frac{90 + 280}{500} = \frac{370}{500} = 74\,\%$。

\textbf{12.} 医院 B：轻症 $\frac{340}{400} = 85\,\%$；
重症 $\frac{65}{100} = 65\,\%$；总体
为 $\frac{340 + 65}{500} = \frac{405}{500} = 81\,\%$。
而医院 A 在轻症（$90 > 85$）上胜出，\emph{并且}在重症
（$70 > 65$）上同样胜出。

\textbf{13.} 总体上，医院 B 为 $81\,\%$，而医院 A 为 $74\,\%$：
在\emph{每一}类中都落败的医院，反而在总体上胜出。
（这种逆转有一个名字：辛普森悖论。）

\textbf{14.} 两家医院的构成恰好相反：医院 A 的患者多为
重症（$400$ 人，占 $500$ 人），医院 B 则多为轻症（$400$ 人，占 $500$ 人）。
无论在哪家医院，重症的治愈率都较低，因此医院 A
的总体数字被其收治的疑难病例拉低——
总体数字是一种加权平均，而其权重在两家
医院之间各不相同。重症患者应选择医院 A：$70\,\%$
在与其相关的唯一一行中胜过 $65\,\%$。

\textbf{15.} 例如：“你的比较混合了轻症和
重症患者，而两家医院的患者构成恰好相反——
按病情分开看，医院 A 在
\emph{两类}患者中的治愈率都更高。总体数字只有在其背后
的群体相似时才可比较；否则决定胜负的是权重，
而非医疗质量。”
\end{solution}
